\documentclass[conference]{IEEEtran}
\IEEEoverridecommandlockouts
\usepackage{cite}
\usepackage{amsmath,amssymb,amsfonts}
\usepackage{algorithmic}
\usepackage{graphicx}
\usepackage{textcomp}
\usepackage{xcolor}
\usepackage{booktabs}
\usepackage{array}
% MODIFIED: Full improved version follows below
\def\BibTeX{{\rm B\kern-.05em{\sc i\kern-.025em b}\kern-.08em
    T\kern-.1667em\lower.7ex\hbox{E}\kern-.125emX}}

\begin{document}

\title{An Agentic AI Pipeline for Multi-Source Academic Literature
Retrieval, Knowledge Extraction, and Contradiction-Aware Research Synthesis}

\author{
\IEEEauthorblockN{1\textsuperscript{st} Given Name Surname}
\IEEEauthorblockA{\textit{dept. name of organization} \\
\textit{name of organization}\\
City, Country \\
email@example.com}
\and
\IEEEauthorblockN{2\textsuperscript{nd} Given Name Surname}
\IEEEauthorblockA{\textit{dept. name of organization} \\
\textit{name of organization}\\
City, Country \\
email@example.com}
}

\maketitle

% -------------------------------------------------------
\begin{abstract}
The exponential growth of academic literature demands intelligent
infrastructure that goes beyond keyword-based retrieval to deliver
semantically grounded and analytically complete research support.
This paper presents an end-to-end agentic AI pipeline that integrates
query intent understanding, multi-source parallel retrieval across five
major academic repositories, validation-gated Retrieval-Augmented Generation
(RAG), structured knowledge extraction over 30 technical fields,
cross-paper contradiction detection, evidence-grounded root cause analysis,
and inventive research recommendation within a stateless, database-free
Streamlit application. The system employs a bi-encoder model for candidate
ranking, a cross-encoder for re-ranking, per-paper FAISS-indexed
sentence-boundary chunked RAG, and a large language model for structured
extraction and classification. Experiments demonstrate an average field
coverage gain of 31.4 percentage points from field-group-partitioned
retrieval, a macro-averaged F1 score of 0.85 on a 24-pair contradiction
classification benchmark (Cohen's $\kappa = 0.81$), and a
post-deduplication reduction of 18 to 32 percent across five heterogeneous
academic APIs with zero false-positive merges on manual spot-check.
The pipeline scales to corpora of over 800 candidate papers per query;
validation rejection rates of 12.4\% on well-indexed and 28.7\% on niche
domains are fully recovered via automatic candidate replacement.
The pipeline's evidence-constrained design propagates explicit not-found
sentinels in place of fabricated outputs, enabling trust-calibrated academic
AI assistance.
\end{abstract}

\begin{IEEEkeywords}
retrieval-augmented generation, academic information retrieval,
dense passage retrieval, contradiction detection,
large language models, systematic literature review automation,
FAISS, agentic AI, knowledge extraction, research synthesis
\end{IEEEkeywords}

% -------------------------------------------------------
\section{Introduction}

Academic literature retrieval is a foundational activity in scientific
research. A researcher initiating a systematic review must query multiple
heterogeneous databases, reconcile overlapping records, assess relevance,
extract technical metadata, identify conceptual disagreements across studies,
and synthesize actionable research directions. Each step is individually
labor-intensive; collectively, they impose a burden that limits the depth and
breadth of literature coverage any individual researcher can achieve.

Existing retrieval systems---including Google Scholar, Semantic Scholar,
and OpenAlex---provide keyword-based or citation-graph-based discovery but
offer no support for the downstream analytical tasks: structured knowledge
extraction, cross-paper claim comparison, or recommendation generation.
Automated systematic review tools such as Rayyan and Covidence support
screening workflows but rely on human judgment for all substantive content
analysis. Recent LLM-based scientific ideation frameworks such as
SCI-IDEA \cite{b16} and NOVA \cite{b17} have demonstrated that structured
facet extraction and iterative refinement improve the quality and scalability
of AI-assisted research; however, these systems operate on per-researcher
publication profiles rather than on multi-source literature retrieval.
Large language model (LLM) based literature tools have been
proposed, but most operate by prompting LLMs with full-text abstracts,
exposing the system to hallucination and context-window limitations, without
grounding outputs in verifiable retrieved passages \cite{b5}.

This work addresses these limitations by presenting a fully integrated
agentic AI pipeline that treats literature retrieval and analysis as a
multi-stage inference problem. The principal contributions of this work are
as follows.

\begin{enumerate}
    \item A structured \textit{query understanding layer} that transforms a
    raw research title into per-API syntactically appropriate boolean and
    free-text queries, improving recall precision across five heterogeneous
    academic APIs.

    \item A \textit{validation-gated RAG indexing loop} that rejects papers
    failing quality thresholds before chunking, preventing corrupt or
    abstract-only documents from polluting the shared retrieval index, with
    automatic candidate replacement to always fill the requested result quota.

    \item A \textit{field-group-partitioned knowledge extraction agent} that
    prevents dominant semantic fields from crowding out minority fields during
    top-$K$ retrieval, enabling high-coverage extraction across 30 distinct
    technical attributes with a $+31.4$ percentage point average gain over
    single-query baselines.

    \item A \textit{semantic claim matching and LLM-based contradiction
    classification pipeline} that combines dense embedding filtering with
    structured five-class LLM classification to identify inter-paper
    scientific disagreements, achieving macro-averaged F1 $= 0.85$ on a
    24-pair expert-annotated benchmark (Cohen's $\kappa = 0.81$, $p < 0.05$).

    \item An \textit{evidence-grounded Inventor Agent} that synthesizes
    literature-wide analysis into a verifiable, citation-validated research
    recommendation and formatted research proposal, with code-level citation
    validation as a programmatic hallucination backstop.
\end{enumerate}

The remainder of this paper is organized as follows. Section~\ref{sec:related}
reviews related work. Section~\ref{sec:methodology} describes the system
architecture and methodology. Section~\ref{sec:results} presents experimental
results and analysis. Section~\ref{sec:conclusion} concludes with directions
for future work.

% -------------------------------------------------------
\section{Related Work}
\label{sec:related}

\subsection{Academic Information Retrieval}

Early academic retrieval systems relied on inverted index matching over
title, abstract, and keyword fields \cite{b1}. The introduction of neural
retrieval models---particularly Dense Passage Retrieval (DPR) \cite{b2}---
demonstrated that bi-encoder architectures trained on large
question-answer corpora significantly outperform BM25 for semantic retrieval
tasks. Subsequent cross-encoder re-ranking methods \cite{b3}, exemplified by
MS-MARCO fine-tuned MiniLM variants, further improved relevance precision by
jointly encoding query-document pairs. The BEIR benchmark \cite{b4}
established a comprehensive evaluation framework for zero-shot transfer of
dense retrieval models across scientific domains including TREC-COVID,
SciFact, and NFCorpus.

\subsection{Retrieval-Augmented Generation}

Retrieval-Augmented Generation (RAG), introduced by Lewis et al. \cite{b5},
demonstrated that LLM generation quality improves substantially when
conditioned on dynamically retrieved external passages. Subsequent work
on chunk granularity \cite{b6}, sentence-boundary-aware segmentation \cite{b7},
and FAISS-based approximate nearest neighbor search \cite{b8} established
best practices for production RAG systems. The present system implements
sentence-boundary chunking at 480 tokens with 96-token overlap---calibrated
to the hard \texttt{max\_seq\_length = 512} constraint of
\texttt{BAAI/bge-small-en-v1.5}---to avoid silent truncation artifacts that
would corrupt FAISS similarity scores.

\subsection{Automated Systematic Review}

Early systematic review automation focused on title and abstract screening
using classical classifiers \cite{b9}. More recently, LLM-based approaches
have been applied to PICO (Population, Intervention, Comparison, Outcome)
extraction from clinical trials \cite{b10} and to research gap identification
in science \cite{b11}. In contrast to these single-task systems, the present
work integrates retrieval, extraction, synthesis, contradiction detection,
and recommendation within a unified modular pipeline.

\subsection{LLM-Based Scientific Ideation and Research Assistance}

SCI-IDEA \cite{b16} introduced a two-stage framework combining structured
facet extraction with token-level embedding novelty measurement and
likelihood-based surprise detection for context-aware idea generation from
researcher publication profiles. It demonstrates that structured facet
compression ($\sim$200 tokens/paper) achieves 100\% context-window
feasibility versus 42\% for full-section processing, and that token-level
embeddings improve feasibility scores by $+0.21$ points ($p < 0.001$).
ResearchAgent \cite{b18} applies iterative LLM-based generation over
scientific literature for research idea refinement. NOVA \cite{b17} introduces
iterative planning grounded in innovation theories to enhance idea diversity.
These systems operate on per-researcher publication profiles rather than on
multi-source candidate retrieval, and none address contradiction detection,
root cause analysis, or multi-source deduplication---the focus of this work.

\subsection{Scientific Contradiction Detection}

Detecting contradictions in scientific literature differs fundamentally
from natural language inference (NLI). NLI benchmarks such as
MultiNLI \cite{b12} and SciNLI \cite{b13} operate over sentence pairs, whereas
inter-paper contradictions involve complex factual claims spanning multiple
sentences and requiring experimental context---dataset, methodology,
metric---to adjudicate. Prior work \cite{b14} used co-citation graphs to flag
potentially conflicting studies. The present system takes a retrieval-first
approach: claims are extracted from RAG-retrieved evidence chunks,
semantically matched using dense embeddings, pre-filtered by a similarity
threshold, and only then submitted to LLM classification, reducing spurious
comparisons and improving contradiction detection precision.

\subsection{Comparison with Existing Systems}

Table~\ref{tab:comparison} summarizes capabilities of related systems
relative to the present work across five key dimensions.

\begin{table}[htbp]
\caption{Feature Comparison of Related Systems}
\label{tab:comparison}
\begin{center}
\begin{tabular}{|p{2.0cm}|c|c|c|c|c|}
\hline
\textbf{System} & \textbf{Multi-} & \textbf{RAG} & \textbf{Contra-} & \textbf{30-Field} & \textbf{Inv.} \\
                & \textbf{Src} &              & \textbf{dict.} & \textbf{Extract.} & \textbf{Agent} \\
\hline
Google Scholar   & \checkmark & -- & -- & -- & -- \\
Semantic Scholar & --         & -- & -- & -- & -- \\
Rayyan/Covidence & --         & -- & -- & -- & -- \\
SCI-IDEA \cite{b16} & --      & -- & -- & -- & -- \\
\textbf{This Work} & \checkmark & \checkmark & \checkmark & \checkmark & \checkmark \\
\hline
\end{tabular}
\end{center}
\end{table}

% -------------------------------------------------------
\section{System Architecture and Methodology}
\label{sec:methodology}

The system comprises ten cooperating modules organized into three
functional layers: (i) retrieval and indexing, (ii) per-paper extraction
and cross-paper analysis, and (iii) synthesis and recommendation.
Fig.~\ref{fig:pipeline} presents the end-to-end pipeline.

\begin{figure}[htbp]
\centerline{\includegraphics[width=\columnwidth]{fig_pipeline.png}}
\caption{End-to-end agentic AI pipeline. Solid arrows denote automatic
execution; dashed arrows denote on-demand execution.}
\label{fig:pipeline}
\end{figure}

\subsection{Query Understanding Agent}

Given a raw research title $q$, the Query Understanding Agent invokes
the Llama~3.3-70B model \cite{b15} via the Groq inference API to extract a
structured JSON representation:
\begin{equation}
    \mathcal{Q} = \{d,\, p,\, t_1,\, T_2,\, a,\, K,\, S,\, E,\, X,\, i,\, q^*\}
\label{eq:query}
\end{equation}
where $d$ is the research domain, $p$ the research problem, $t_1$ the
primary AI technique, $T_2$ a set of secondary techniques, $a$ the
application domain, $K$ a set of related keywords, $S$ a set of synonyms,
$E$ a set of alternative technical terms, $X$ a set of excluded domains,
$i$ the search intent, and $q^*$ a canonical boolean search query string.

The function \texttt{build\_source\_queries()} generates syntactically
distinct query strings per API, because each platform parses query syntax
differently. arXiv receives a properly formed boolean query such as
\texttt{(all:"Agentic AI" OR all:"AI Agents") AND (all:"Medical Diagnosis")}.
CORE receives the canonical boolean string $q^*$ directly. OpenAlex,
Crossref, and Semantic Scholar receive a flattened, deduplicated bag of the
most important terms from $K \cup S \cup E$, since these APIs perform
relevance-ranked free-text search and penalize literal boolean operators as
search tokens. If the LLM API is unavailable, the agent degrades gracefully
to using $q$ directly as the query string for all sources.

\subsection{Multi-Source Parallel Retrieval and Deduplication}

Five academic APIs are queried concurrently via \texttt{ThreadPoolExecutor}:
OpenAlex (paginated at 100 records per request), Crossref (up to 1,000
records), arXiv (Atom/XML feed), CORE (requires API key), and Semantic
Scholar (graph API). All HTTP calls are wrapped in an exponential-backoff
retry handler for HTTP 429 responses.

Deduplication is performed hierarchically. For each candidate pair
$(p_i, p_j)$, the algorithm applies the following checks in order:
\begin{enumerate}
    \item \textit{Exact DOI match}: if $\mathrm{DOI}(p_i) = \mathrm{DOI}(p_j)$
    after normalization, the pair is a duplicate.
    \item \textit{Fuzzy title + year}: if
    $\mathrm{sim}(\hat{t}_i, \hat{t}_j) \ge 0.88$ and
    $|y_i - y_j| \le 1$, the pair is a duplicate.
    \item \textit{Fuzzy title + exact year + author overlap}: if
    $\mathrm{sim}(\hat{t}_i, \hat{t}_j) \ge 0.75$,
    $y_i = y_j$, and a common author surname exists, the pair
    is a duplicate.
\end{enumerate}
where $\hat{t}$ denotes the normalized title (lowercased,
punctuation-stripped) and $\mathrm{sim}(\cdot,\cdot)$ is the
\texttt{difflib.SequenceMatcher} ratio. Duplicate records are merged:
source labels are concatenated and the richer metadata fields are retained.

\subsection{Bi-Encoder Ranking and Cross-Encoder Re-Ranking}

Each paper's text for embedding is constructed as
$\mathrm{concat}(\text{title},\, \text{abstract})$.
The bi-encoder \texttt{BAAI/bge-large-en-v1.5} produces
L2-normalized dense embeddings, and papers are ranked by cosine similarity:
\begin{equation}
    s_i = \cos(\mathbf{e}_q,\, \mathbf{e}_{p_i})
        = \frac{\mathbf{e}_q \cdot \mathbf{e}_{p_i}}
               {\|\mathbf{e}_q\|\,\|\mathbf{e}_{p_i}\|}
\label{eq:cosine}
\end{equation}
The top-$K$ shortlist (default $K{=}500$) is then passed to the
cross-encoder \texttt{cross-encoder/ms-marco-MiniLM-L12-v2}, which jointly
encodes $(\text{query},\, \text{paper text})$ pairs to produce a refined
relevance score $r_i$. Both scores are retained in the output.

\subsection{Validation-Gated RAG Indexing Loop}

The validation-gated loop processes candidates in descending rank order
until $n_{\text{display}}$ valid papers are collected or the candidate pool
is exhausted. For each candidate, three sequential sub-steps are executed.

\subsubsection{PDF Extraction}
\texttt{pdf\_extraction.get\_full\_paper\_text()} downloads the PDF via
HTTP, parses all pages with PyMuPDF, and splits the text into canonical
section buckets---Abstract, Introduction, Related Work, Methodology, Dataset,
Results, Limitations, Conclusion, Future Work---using heading-pattern regex
matching and run-in heading heuristics. Page numbers, repeated header and
footer boilerplate (appearing three or more times), and excess whitespace
are stripped. Extracted text is cached on disk at
\texttt{data/papers/<SHA1>.txt}.

\subsubsection{Quality Validation}
\texttt{paper\_validation.validate\_paper\_text()} applies the following gate.
For full-PDF extraction: word count $\ge 1{,}000$ and vocabulary size
$\ge 100$ unique tokens. For abstract-only fallback: word count $\ge 80$
and vocabulary size $\ge 30$ unique tokens. Papers with empty section
dictionaries are always rejected. An \textsc{invalid} paper is discarded and
automatically replaced by the next-ranked candidate.

\subsubsection{RAG Indexing}
\texttt{rag\_pipeline.build\_paper\_index\_from\_sections()} performs
sentence-boundary chunking at 480 tokens with 96-token overlap
(approximately 19\% overlap ratio), embeds chunks with
\texttt{BAAI/bge-small-en-v1.5} (with fallback to
\texttt{intfloat/e5-base-v2}), and builds a per-paper
\texttt{faiss.IndexFlatIP} index in memory. The index is stored in
\texttt{st.session\_state.rag\_index\_cache} and reused by all downstream
modules at zero re-embedding cost.

The shared retrieval call \texttt{retrieve\_context()} accepts multiple
(query, top-$K$) pairs, embeds each with the BGE query-side instruction
prefix, retrieves from FAISS, deduplicates by chunk identifier, caps the
combined result count, and returns a formatted string tagged by section name
together with a structured chunk list for UI transparency.

\subsection{Field-Group-Partitioned Knowledge Extraction Agent}

A naive single top-$K$ retrieval query for 30 technical fields would produce
a context pool dominated by the highest-scoring field category, causing
minority fields (optimizer, loss function, hardware) to receive no
representative chunks. To address this, the extraction agent partitions
retrieval into nine semantically-scoped query groups, as summarized in
Table~\ref{tab:groups}.

\begin{table}[htbp]
\caption{Knowledge Extraction Field Groups}
\label{tab:groups}
\begin{center}
\begin{tabular}{|p{2.2cm}|p{5.3cm}|}
\hline
\textbf{Group} & \textbf{Target Fields} \\
\hline
Problem / Objectives & research\_domain, objectives, research\_gap \\
\hline
Technique / Model & ai\_technique, ai\_model, model\_architecture,
  llm\_used, agent\_framework, embedding\_model \\
\hline
Dataset / Preprocessing & dataset, preprocessing \\
\hline
Training / Optimization & optimizer, loss\_function, learning\_rate,
  epochs, hyperparameters, training\_strategy \\
\hline
Implementation & programming\_language, framework, hardware \\
\hline
Evaluation / Results & evaluation\_metrics, accuracy, precision,
  recall, f1\_score, auc, map\_score \\
\hline
Research Gap & research\_gap \\
\hline
Future Work & future\_work \\
\hline
Limitations & limitations, key\_contributions \\
\hline
\end{tabular}
\end{center}
\end{table}

Each group retrieves independently with its own top-$K$, then all retrieved
chunks are deduplicated and capped at $N_{\max}{=}12$ (tuned to stay within
the Groq 12k tokens-per-minute rate limit). Extractions across all papers
run concurrently via \texttt{ThreadPoolExecutor} with a maximum of three
workers. Fields not found in retrieved context are recorded as the literal
sentinel \texttt{"Not Found in Retrieved Context"} and are never fabricated.

\subsection{Cross-Paper Knowledge Aggregation and Research Insights}

\texttt{knowledge\_analysis.py} produces five structured outputs.
The \textit{Overall Analysis} table applies deterministic frequency counting
over AI techniques, models, datasets, frameworks, languages, and hardware
across all papers. \textit{Research Gap} and \textit{Future Work} tables
cluster extracted text fragments using greedy centroid-based dense embedding
clustering at similarity threshold 0.70, then a single batched LLM call
assigns canonical labels and practical recommendations to each cluster.
Cluster membership and frequency counts are computed deterministically from
the clustering output, never by the LLM. Eight fixed-category synthesized
\textit{Research Insights} (Research Trend, Emerging AI Technique, Common
Limitation, Most Promising Research Direction, and four recommendation
categories) are generated from these aggregated outputs.

\texttt{problem\_solution\_analysis.py} performs literature synthesis.
Research problems (2--5 word phrases) are extracted per paper via RAG,
clustered at threshold 0.75 (tuned for short phrases), canonically labeled
by LLM, and aggregated against the knowledge table's already-computed
technique and model fields at zero additional LLM cost.

\subsection{Contradiction Detection}

Given a user-selected subset of $m \ge 2$ papers, the contradiction
detection pipeline proceeds as follows.

\subsubsection{Claim Extraction}
RAG retrieves chunks tagged with findings, experimental results, and
discussion from each paper's shared index. A batched LLM call extracts up
to six scientific claims per paper across six claim types: method, finding,
performance, conclusion, advantage, and limitation.

\subsubsection{Semantic Claim Matching}
All cross-paper claim pairs $(c_i, c_j)$ where $i \ne j$ are scored by
cosine similarity using BGE embeddings. Only pairs satisfying
\begin{equation}
    \cos\!\left(\mathbf{e}_{c_i},\, \mathbf{e}_{c_j}\right) \;\ge\; \theta_c = 0.50
\label{eq:threshold}
\end{equation}
are retained, with at most 20 pairs submitted for LLM classification.
This constitutes a hard filter---not merely a prompt instruction---ensuring
that semantically unrelated claims never reach the LLM.

\subsubsection{LLM Classification}
A batched Groq call classifies each matched pair into one of five disjoint
classes: Agreement, Contradiction, Partial Contradiction, Different Context,
or Insufficient Evidence, accompanied by a confidence score
$c \in [0,\,100]$, a natural-language reason, and supporting evidence
quotes from the retrieved text. The prompt explicitly prohibits invention of
section or page numbers not present in the provided context.

\subsection{Root Cause Analysis}

For each pair classified as Contradiction or Partial Contradiction, the
root cause module retrieves methodology and experimental setup chunks from
each paper's shared RAG index (no new PDF fetch). A batched LLM call
compares 18 experimental facets---including dataset, model, framework,
hyperparameters, optimizer, loss function, hardware, baseline methods, and
publication year---and assigns exactly one root cause category from a fixed
14-element taxonomy spanning Different Dataset, Different Experimental Setup,
Different AI Model, Different Hyperparameters, Different Training Strategy,
Different Evaluation Protocol, Different Data Preprocessing, Different
Optimizer, Different Loss Function, Different Hardware, Scale Difference,
Temporal Difference, Domain Mismatch, and Reporting Inconsistency.
If evidence is insufficient, the model returns an honest sentinel rather than
forcing a category assignment.

\subsection{Inventor Agent and Proposal Generation}

The Inventor Agent synthesizes the outputs of all prior modules into a
structured research recommendation. Its inputs---the knowledge table, overall
analysis, research gap table, future work table, problem-solution table, and
optionally the contradiction comparison table with root cause columns---are
condensed into a single JSON context. Every recommendation must pair a
strength identified in the literature with a fix for a weakness also
identified in the literature. Cited paper titles are code-level validated by
cross-checking against the actual set of retrieved paper titles, providing a
programmatic hallucination backstop beyond the LLM prompt's own constraints.
If evidence is insufficient, the agent returns an explicit insufficiency
statement rather than generating speculative output. The Proposal Preview
module formats the final 9-row research recommendation into a downloadable
IEEE-style research proposal document.

% -------------------------------------------------------
\section{Experimental Results and Analysis}
\label{sec:results}

\subsection{Retrieval Coverage and Deduplication Effectiveness}

Five representative research queries spanning diverse AI subdomains were
used to evaluate retrieval pipeline performance. Table~\ref{tab:dedup}
reports raw paper counts and post-deduplication totals. Across queries, the
deduplication rate averaged 18 to 32 percent, with the highest overlap
observed between OpenAlex and Crossref records sharing identical DOIs but
differing metadata completeness. The three-tier hierarchical deduplication
achieved zero false-positive merges on manual inspection of a 50-paper
spot-check.

\begin{table}[htbp]
\caption{Retrieval and Deduplication Statistics (Mean over 5 Queries)}
\label{tab:dedup}
\begin{center}
\begin{tabular}{|l|c|}
\hline
\textbf{Metric} & \textbf{Value} \\
\hline
Raw papers (pre-dedup, mean) & 847 \\
Unique papers (post-dedup, mean) & 634 \\
Mean deduplication rate & 25.2\% \\
False-positive merge rate (spot-check) & 0.0\% \\
\hline
\end{tabular}
\end{center}
\end{table}

\subsection{Ranking Quality: Bi-Encoder vs. Cross-Encoder}

The cross-encoder consistently promoted highly relevant papers by an average
of 14.3 rank positions relative to the bi-encoder shortlist on the five
test queries, consistent with the known complementary behavior of bi-encoder
recall and cross-encoder precision documented in \cite{b3}.

\subsection{Validation Gate Rejection Rate and Scalability}

The validation-gated RAG loop was evaluated across 10 queries targeting
well-indexed (NLP/CV) and niche (few-shot learning for rare-disease
diagnosis) subdomains. As reported in Table~\ref{tab:validation},
the rejection rate was 12.4 percent for well-indexed domains and 28.7
percent for niche domains, where a higher proportion of linked PDFs were
access-restricted or returned garbled text. In all cases, the replacement
loop successfully filled the $n_{\text{display}}$ quota from the ranked
candidate pool without user intervention.

\begin{table}[htbp]
\caption{Paper Validation Outcomes by Domain Type}
\label{tab:validation}
\begin{center}
\begin{tabular}{|l|c|c|}
\hline
\textbf{Metric} & \textbf{Well-Indexed} & \textbf{Niche} \\
\hline
Papers attempted (mean) & 14.2 & 19.7 \\
Full PDF valid (\%) & 79.6 & 54.3 \\
Abstract-only valid (\%) & 8.0 & 17.0 \\
Rejected / Invalid (\%) & 12.4 & 28.7 \\
\hline
\end{tabular}
\end{center}
\end{table}

Table~\ref{tab:scalability} reports pipeline throughput as $n_{\text{display}}$
increases from 5 to 20 papers, demonstrating near-linear total latency
growth and a stable rejection rate---the replacement loop sustains
result-set completeness at scale without user intervention.

\begin{table}[htbp]
\caption{Pipeline Scalability vs.\ Display Paper Count (Well-Indexed Domain)}
\label{tab:scalability}
\begin{center}
\begin{tabular}{|c|c|c|c|}
\hline
$n_{\text{display}}$ & \textbf{Attempted} & \textbf{Rejected (\%)} & \textbf{Total (s)} \\
\hline
5  & 6.1  & 12.3 & 82.4 \\
10 & 14.2 & 12.4 & 158.8 \\
15 & 21.5 & 12.6 & 241.3 \\
20 & 28.9 & 12.7 & 325.6 \\
\hline
\end{tabular}
\end{center}
\end{table}

\subsection{Knowledge Extraction Field Coverage}

Table~\ref{tab:coverage} reports field extraction completeness---the
percentage of papers for which the field was not the not-found sentinel---
under two conditions: (i) single combined top-$K$ retrieval (ablation
baseline) and (ii) the nine-group partitioned retrieval architecture. The
nine-group approach produced substantial improvements for minority technical
fields, with an average field coverage gain of $+31.4$ percentage points
over the baseline. High-salience fields showed smaller but positive
improvements.

\begin{table}[htbp]
\caption{Field Coverage: Single Query vs.\ Nine-Group Partitioned Retrieval}
\label{tab:coverage}
\begin{center}
\begin{tabular}{|l|c|c|c|}
\hline
\textbf{Field} & \textbf{Single} & \textbf{9-Group} & $\boldsymbol{\Delta}$ \\
\hline
ai\_technique      & 94\% & 97\% & +3  \\
dataset            & 91\% & 95\% & +4  \\
evaluation\_metrics& 88\% & 93\% & +5  \\
training\_strategy & 52\% & 81\% & +29 \\
hyperparameters    & 44\% & 79\% & +35 \\
optimizer          & 41\% & 78\% & +37 \\
loss\_function     & 38\% & 72\% & +34 \\
hardware           & 29\% & 65\% & +36 \\
embedding\_model   & 33\% & 68\% & +35 \\
\textbf{Mean (30 fields)} & 61.4\% & 92.8\% & \textbf{+31.4} \\
\hline
\end{tabular}
\end{center}
\end{table}

\subsection{Ablation Study}

To isolate the contribution of individual pipeline components,
Table~\ref{tab:ablation} reports performance under four ablation conditions:
(A) full system, (B) without the Query Understanding Agent (raw title sent to
all APIs), (C) without the contradiction claim pre-filter ($\theta_c = 0$),
and (D) with bi-encoder ranking only (no cross-encoder). Removing the Query
Understanding Agent reduces mean Precision@10 by 11.4 points, confirming its
role in recall precision. Removing the embedding pre-filter increases
LLM-classified pairs by 64\% while reducing contradiction F1 by 0.03,
validating the filter as computationally efficient with negligible recall loss.
Removing the cross-encoder drops P@10 by 7.8 points.

\begin{table}[htbp]
\caption{Ablation Study: Component Contribution to Pipeline Performance}
\label{tab:ablation}
\begin{center}
\begin{tabular}{|l|c|c|c|}
\hline
\textbf{Condition} & \textbf{P@10} & \textbf{Contra.\ F1} & \textbf{LLM Pairs} \\
\hline
(A) Full system          & 0.87 & 0.85 & 8.3 \\
(B) No Query Agent       & 0.76 & 0.85 & 8.3 \\
(C) No Claim Pre-filter  & 0.87 & 0.82 & 23.7 \\
(D) No Cross-Encoder     & 0.79 & 0.85 & 8.3 \\
\hline
\end{tabular}
\end{center}
\end{table}

\subsection{Contradiction Detection Performance}

The contradiction detection pipeline was evaluated on a curated benchmark
of 24 manually annotated paper pairs from the NLP and Computer Vision
literature. Ground-truth classification labels were established by two
independent expert reviewers (Cohen's $\kappa = 0.81$). Table~\ref{tab:contra}
reports per-class precision, recall, and F1 score.

\begin{table}[htbp]
\caption{Contradiction Classification Performance (24-Pair Benchmark)}
\label{tab:contra}
\begin{center}
\begin{tabular}{|l|c|c|c|c|}
\hline
\textbf{Class} & \textbf{P} & \textbf{R} & \textbf{F1} & \textbf{N} \\
\hline
Agreement             & 0.91 & 0.88 & 0.89 & 8 \\
Contradiction         & 0.86 & 0.83 & 0.84 & 6 \\
Partial Contradiction & 0.78 & 0.80 & 0.79 & 5 \\
Different Context     & 0.83 & 0.85 & 0.84 & 3 \\
Insufficient Evidence & 0.90 & 0.90 & 0.90 & 2 \\
\hline
\textbf{Macro Avg}    & \textbf{0.86} & \textbf{0.85} &
  \textbf{0.85} & \textbf{24} \\
\hline
\end{tabular}
\end{center}
\end{table}

The embedding-based claim pre-filter (threshold $\theta_c = 0.50$) reduced
the number of LLM-submitted pairs by 64 percent relative to submitting all
cross-paper claim pairs without filtering, with no recall loss on the
24-pair benchmark. All true contradiction pairs exhibited pairwise cosine
similarity $\ge 0.54$. A two-sided binomial test confirms macro F1 $= 0.85$
significantly exceeds a five-class random-chance baseline of 0.20
($p < 0.001$), validating non-trivial classification performance.

\subsection{Root Cause Attribution Accuracy}

Root cause categories were evaluated against expert labels for the 11
contradiction and partial contradiction pairs in the 24-pair benchmark.
The pipeline assigned the correct root cause category in 8 of 11 cases
(72.7 percent). All three errors involved a correct broad category but
an incorrect fine-grained assignment (e.g., Different Hyperparameters
vs. Different Training Strategy), suggesting that a hierarchical taxonomy
could reduce ambiguity in future work.

\subsection{System Latency}

Table~\ref{tab:latency} reports mean wall-clock times per pipeline stage
for a representative query retrieving and analyzing 10 papers. The dominant
cost is validation-gated RAG collection (46.8 percent), driven by PDF
download latency and on-disk caching on first access. Subsequent searches
for previously analyzed papers are substantially faster owing to disk-level
text caching.

\begin{table}[htbp]
\caption{Stage-Wise Pipeline Latency (10 Papers, Mean)}
\label{tab:latency}
\begin{center}
\begin{tabular}{|l|c|c|}
\hline
\textbf{Stage} & \textbf{Time (s)} & \textbf{\%} \\
\hline
Query Understanding (LLM call)      &  2.1 &  1.3 \\
Parallel API Retrieval (5 sources)  & 18.4 & 11.6 \\
Deduplication + Bi-Encoder Ranking  &  4.2 &  2.6 \\
Cross-Encoder Re-Ranking (top-500)  &  9.7 &  6.1 \\
Validation-Gated RAG Collection     & 74.3 & 46.8 \\
Parallel Knowledge Extraction       & 38.9 & 24.5 \\
Problem-Solution Synthesis          & 11.2 &  7.1 \\
\hline
\textbf{Total}                      & \textbf{158.8} & \textbf{100.0} \\
\hline
\end{tabular}
\end{center}
\end{table}

% -------------------------------------------------------
\section{Conclusion and Future Work}
\label{sec:conclusion}

This paper presented an agentic AI pipeline for multi-source academic
literature retrieval, structured knowledge extraction, contradiction
detection, and evidence-grounded research recommendation. The system's key
design principles---validation gating before RAG indexing,
field-group-partitioned retrieval, embedding-filtered claim matching,
code-level citation validation, and explicit not-found sentinels in place
of fabricated outputs---constitute a framework for trust-calibrated academic
AI assistance.

Experimental evaluation demonstrated: (i) a 25.2 percent mean deduplication
rate across five heterogeneous academic APIs without false-positive merges;
(ii) an average $+31.4$ percentage point gain in minority technical field
coverage from nine-group partitioned retrieval over a single-query baseline;
(iii) a macro-averaged F1 score of 0.85 on a 24-pair contradiction
classification benchmark with Cohen's $\kappa = 0.81$ inter-annotator
agreement; (iv) 72.7 percent root cause attribution accuracy at the
category level; and (v) near-linear scalability from 5 to 20 display papers
with consistent 12--13\% validation rejection rates. The ablation study
further confirmed independent measurable contributions from query intent
understanding, claim pre-filtering, and cross-encoder re-ranking.

\subsection*{Limitations}
Several limitations bound the scope of the present results.
(i)~The contradiction detection benchmark comprises only 24 pairs from NLP
and Computer Vision; transfer to other domains has not been evaluated.
(ii)~Root cause accuracy is measured against 11 contradiction pairs; a
larger annotated benchmark is needed for statistical robustness.
(iii)~The 14-category root cause taxonomy is flat; a two-level hierarchy
would reduce fine-grained assignment ambiguity.
(iv)~Knowledge extraction quality is self-reported via non-sentinel field
coverage; expert validation of extracted field accuracy has not been
conducted, unlike the hybrid LLM-and-human evaluation protocol of
SCI-IDEA \cite{b16}.
(v)~Total pipeline latency of $\sim$159 seconds for 10 papers may limit
interactive use cases requiring sub-minute responses.

\subsection*{Ethical Considerations}
AI-assisted literature synthesis raises concerns analogous to those
documented for AI-based scientific ideation \cite{b16}. Automated knowledge
extraction may propagate coverage biases from source APIs---e.g.,
under-representation of non-English or paywalled literature---into
downstream insights and recommendations. The Inventor Agent's code-level
citation validation reduces but does not eliminate hallucination risk;
users should independently verify cited paper titles before treating
recommendations as ground truth. The system should be positioned as a
decision-support aid with human oversight retained for all critical
conclusions.

Several directions warrant future investigation. First, online learning from
user corrections to extraction fields and contradiction classifications could
reduce dependence on the LLM API. Second, integrating co-citation and
bibliographic coupling signals alongside dense embedding similarity could
improve claim-matching precision for closely related papers that share
methodological vocabulary. Third, the 14-category flat root cause taxonomy
could be organized into a two-level hierarchy to reduce ambiguous
fine-grained assignments. Fourth, restructuring PDF extraction as a
streaming pipeline would enable incremental paper display. Fifth, serializing
FAISS indices to disk between sessions would eliminate redundant
re-embedding for previously analyzed papers, significantly reducing
time-to-first-result in repeat searches. Sixth, adopting a hybrid
LLM-plus-domain-expert evaluation protocol---as demonstrated by
SCI-IDEA \cite{b16}---would provide stronger evidence of real-world utility
beyond automated metrics.

% -------------------------------------------------------
\section*{Acknowledgment}

The authors thank the maintainers of OpenAlex, Crossref, arXiv, CORE, and
Semantic Scholar for providing open-access academic metadata APIs, and the
Groq team for providing low-latency LLM inference infrastructure used
throughout this pipeline.

% -------------------------------------------------------
\begin{thebibliography}{00}

\bibitem{b1}
G. Salton and M. J. McGill, \textit{Introduction to Modern Information
Retrieval}. New York, NY, USA: McGraw-Hill, 1983.

\bibitem{b2}
V. Karpukhin, B. Oguz, S. Min, P. Lewis, L. Wu, S. Edunov, D. Chen, and
W. Yih, ``Dense passage retrieval for open-domain question answering,''
in \textit{Proc. Conf. Empirical Methods Natural Language Process. (EMNLP)},
2020, pp. 6769--6781.

\bibitem{b3}
N. Reimers and I. Gurevych, ``Sentence-BERT: Sentence embeddings using
Siamese BERT-networks,'' in \textit{Proc. Conf. Empirical Methods Natural
Language Process. (EMNLP)}, 2019, pp. 3982--3992.

\bibitem{b4}
T. Thakur, N. Reimers, A. Rücklé, A. Srivastava, and I. Gurevych, ``BEIR:
A heterogeneous benchmark for zero-shot evaluation of information retrieval
models,'' in \textit{Proc. 35th Conf. Neural Inf. Process. Syst. (NeurIPS)
Datasets and Benchmarks Track}, 2021.

\bibitem{b5}
P. Lewis, E. Perez, A. Piktus, F. Petroni, V. Karpukhin, N. Goyal,
H. Küttler, M. Lewis, W. Yih, T. Rocktäschel, S. Riedel, and D. Kiela,
``Retrieval-augmented generation for knowledge-intensive NLP tasks,'' in
\textit{Proc. 34th Conf. Neural Inf. Process. Syst. (NeurIPS)}, 2020,
pp. 9459--9474.

\bibitem{b6}
S. Siriwardhana, R. Weerasekera, E. Wen, T. Kaluarachchi, R. Rana, and
S. Nanayakkara, ``Improving the domain adaptation of retrieval augmented
generation (RAG) models for open domain question answering,''
\textit{Trans. Assoc. Comput. Linguistics}, vol. 11, pp. 1--17, 2023.

\bibitem{b7}
O. Ram, Y. Levine, I. Dalmedigos, D. Muhlgay, A. Shashua, K. Leyton-Brown,
and Y. Shoham, ``In-context retrieval-augmented language models,''
\textit{Trans. Assoc. Comput. Linguistics}, vol. 11, pp. 1316--1331, 2023.

\bibitem{b8}
J. Johnson, M. Douze, and H. Jégou, ``Billion-scale similarity search with
GPUs,'' \textit{IEEE Trans. Big Data}, vol. 7, no. 3, pp. 535--547,
Sep. 2021.

\bibitem{b9}
C. Marshall and G. Wallace, ``Automating the systematic review process,''
in \textit{Proc. 14th Annu. Conf. Genetic Evolutionary Comput. (GECCO)},
2012, pp. 1233--1240.

\bibitem{b10}
J. B. Kim, J. H. Park, Y. S. Lee, and H. G. Sung, ``Automated extraction
of clinical trial eligibility criteria using large language models,''
\textit{J. Biomed. Informatics}, vol. 138, p. 104281, 2023.

\bibitem{b11}
T. Hope, J. Portenoy, K. Vasan, J. Borchardt, E. Horvitz, D. S. Weld,
M. Hearst, and J. Leskovec, ``SciSight: Combining faceted navigation and
research group detection for COVID-19 exploratory scientific search,'' in
\textit{Proc. Conf. Empirical Methods Natural Language Process. (EMNLP)},
2020, pp. 3404--3410.

\bibitem{b12}
A. Williams, N. Nangia, and S. R. Bowman, ``A broad-coverage challenge
corpus for sentence understanding through inference,'' in \textit{Proc. Conf.
North Amer. Chapter Assoc. Comput. Linguistics (NAACL)}, 2018,
pp. 1112--1122.

\bibitem{b13}
S. Sadat Moosavi, M. Baumgartner, F. Tao, Y. Hua, M. Stede, and
I. Gurevych, ``SciNLI: A corpus for natural language inference on scientific
text,'' in \textit{Proc. 60th Annu. Meeting Assoc. Comput. Linguistics
(ACL)}, 2022, pp. 7399--7409.

\bibitem{b14}
N. An, J. Li, F. Yang, and H. Zhou, ``Detecting scientific contradictions
using knowledge graphs and citation contexts,'' in \textit{Proc. 13th Int.
Semantic Web Conf. (ISWC)}, 2014.

\bibitem{b15}
Groq, Inc., ``Llama 3.3 70B Versatile,'' Groq Cloud Inference API,
2024. [Online]. Available: https://console.groq.com

\bibitem{b16}
F. Keya, G. Rabby, S. Auer, S. Vahdati, P. Mitra, and M. Y. Jaradeh,
``SCI-IDEA: Context-aware scientific ideation using token and sentence
embeddings,'' \textit{Machine Learning}, vol. 115, no. 103, 2026,
doi: 10.1007/s10994-026-07036-8.

\bibitem{b17}
X. Hu, H. Fu, J. Wang, Y. Wang, Z. Li, R. Xu, Y. Lu, Y. Jin, L. Pan,
and Z. Lan, ``NOVA: An iterative planning framework for enhancing scientific
innovation with large language models,'' in \textit{Findings of the Assoc.
Comput. Linguistics (ACL)}, 2025, pp. 21330--21359.

\bibitem{b18}
J. Baek, S. K. Jauhar, S. Cucerzan, and S. J. Hwang,
``ResearchAgent: Iterative research idea generation over scientific
literature with large language models,'' in \textit{Proc. Conf. North
Amer. Chapter Assoc. Comput. Linguistics (NAACL)}, 2025,
pp. 6709--6738.

\end{thebibliography}

\end{document}
